\section{正交集算法的复杂度}

\begin{theorem}{正交集算法的复杂度}{}
	设$\boldsymbol{U}\in\mathbf{M}_{m,n}\left(\R\right),\boldsymbol{u}_{n+1}\in\mathbf{M}_{n+1,1}\left(\R\right)$.对$\tilde{\boldsymbol{U}}\coloneq\left[\boldsymbol{U}\middle|-\boldsymbol{u}_{n+1}\right]$应用算法(\ref{algo:orthogonal-set-algorithm-for-solving-linear-systems}).如果算法能进行$m$次,则
	\begin{enumerate}
		\item 算法总共需要做$\frac{1}{6}m\left(1+3m-4m^{2}+6mn\right)$次乘法和$\frac{1}{6}m\left(m-1\right)\left(6n-4m-1\right)$次加法.
		\item 算法总共需要$\frac{1}{3}m\left(6mn-4m^{2}+3m-3n+1\right)$次运算.
	\end{enumerate}
\end{theorem}

\begin{remark}{无除法运算的正交集算法}{}
	如果算法(\ref{algo:orthogonal-set-algorithm-for-solving-linear-systems})要求不能使用除法,那么我们约定:对任一$1\leq i\leq m$和$1\leq j\leq N_{i-1}$,规定
	\begin{align*}
		\boldsymbol{w}_{j}^{\left(i\right)}=
		\begin{cases}
			\boldsymbol{w}_{j}^{\left(i-1\right)},&j=r\\
			t_{r}^{\left(i\right)}\boldsymbol{w}_{j}^{\left(i-1\right)}-t_{j}^{\left(i\right)}\boldsymbol{w}_{r}^{\left(i-1\right)},&j\neq r
		\end{cases}
	\end{align*}
\end{remark}